\section{A geometry that stays fixed during the iteration}
\label{sec:fixed-geometry}
\label{sec:rosen-physical-grading}
\label{sec:optical-graph-height}
\label{sec:pp-positive-exact-poisson}

\subsection{Overview of Part~\ref{part:nonlinear}}
\label{subsec:part3-overview}

We now solve the nonlinear equation $\mathcal F_q(w,u)=0$ near the
reference pair $(w_*,\Phi_*)$, where
\[
 \Phi_*=\phi_*\,e_2=\phi_*\sqrt2\cos2\theta,\qquad
 \mathscr L_{w_*}\Phi_*=0,\qquad \ell_*\Phi_*=1 .
\]
The inverse of $\mathscr L_{w_*}$ constructed in
Part~\ref{part:reference} is useful only if the linearization at a
nearby state can be compared with it. Two features of the
linearization would normally move with the state: the circles on which
it changes type, and the angular positions at which characteristics
arrive after crossing a negative component. The first was fixed in
Lemma~\ref{lem:type-completion} by a correction proportional to the
residual. The second is fixed in the same spirit, at the price of
enlarging the problem.

\paragraph{Unknowns.}
The profile is allowed to vary in a controlled family,
\begin{equation}\label{eq:full-height-metric-controls}
 w=w_*\exp\Bigl(\nu\chi_*+\sum_m\kappa_m\chi_m\Bigr),
\end{equation}
where $\chi_*$ is the function of
Proposition~\ref{prop:rad-fixed-theorem}(iii) and, for every negative
component $m$, $\chi_m$ is a bump supported inside that component
(Lemma~\ref{lem:quantitative-optical-controls}). In addition every
negative component carries a circle diffeomorphism $h_{L,m}$, a
\emph{chart} of its left boundary circle. The \emph{state} is
$\bx=(u,\nu,\kappa,h_L)$ with $\kappa=(\kappa_m)_m$ and
$h_L=(h_{L,m})_m$; see Table~\ref{tab:unknowns}.

\begin{table}[htbp]
\centering
\begin{tabularx}{\linewidth}{@{}l l Y@{}}
\toprule
Unknown & Lives in & Role\\
\midrule
$u$ & $\Phi_*+\ker\ell_*\subset\mathcal X_\beta$ &
 The height is $z=h+qu$ with $h'=w$.\\
$\nu$ & $\R$ &
 Amplitude of the metric control $\chi_*$; removes the obstruction in
 the second angular mode.\\
$\kappa=(\kappa_m)$ & $\mathcal K_\beta$ &
 Amplitudes of the metric controls $\chi_m$; adjust the mean travel
 angle of each negative component.\\
$h_L=(h_{L,m})$ & $\id+\mathcal Z_{\beta,\mathrm{odd}}$ &
 Charts of the left boundary circles; straighten the return maps.\\
\bottomrule
\end{tabularx}
\caption{The unknowns of the nonlinear problem. The spaces are defined
in Definition~\ref{def:readable-physical-grading}; $q$ is a parameter.}
\label{tab:unknowns}
\end{table}

\paragraph{Equations.}
There are two. The first is the height equation
$\mathcal A:=\mathcal F_q(w,u)=0$. The second, $\mathfrak e=0$, is
imposed by us: it says that on every negative component the two
characteristic return maps of the linearization, read in the charts,
are the rotations by $\pm T_m$, where $T_m$ is the travel angle of the
reference profile. The \emph{return residual} $\mathfrak e_m$ is an
even function on the circle. Its mean is controlled by the number
$\kappa_m$, and its oscillating part by the odd function
$h_{L,m}-\id$; solving for the latter involves the small divisors
$\sin(jT_m)$ of Proposition~\ref{prop:rad-fixed-theorem}(ii).

\paragraph{The corrected linearization.}
The derivative of $\mathcal A$ with respect to $u$ will be written as
\[
 D_u\mathcal A=\widehat L_{\bx}
 +(\text{a term proportional to }\mathcal A)
 +(\text{a term proportional to }\mathfrak e),
\]
where $\widehat L_{\bx}$ changes type exactly at the turning circles of
$w_*$ and has exactly the return maps $R_{\pm T_m}$ in the charts.
After a change of coordinates all operators $\widehat L_{\bx}$ act on
one set of cylinders with one set of transmission conditions,
independent of $\bx$. This makes the comparison with the reference
operator possible.

\paragraph{Plan.}
This section introduces the weighted spaces
(Section~\ref{subsec:spaces}), the controls and the neighbourhood of
the reference state on which everything is defined
(Sections~\ref{subsec:root-geometry}--\ref{subsec:neighbourhood}), the
return residual and the operator $\widehat L$
(Section~\ref{subsec:return-correction}), and the coordinates
(Sections~\ref{subsec:root-coordinate}--\ref{subsec:positive-crossing}).
Section~\ref{sec:stable-inverse} inverts $\widehat L_{\bx}$,
Section~\ref{sec:solving-the-equation} solves the coupled system by the
Nash--Moser method, and Section~\ref{sec:final-proof} concludes.

The construction involves infinitely many components, and one
smallness condition has to serve for all of them. The room for this
comes from a simple fact, used again and again: the deformation
decays like $H^\beta$, which is far smaller than any power of the
coefficient allowance, while the reference profile decays only like
$e^{-n^{5/2}}$.

\subsection{Weights and function spaces}
\label{subsec:spaces}

\paragraph{Continuous weights.}
Let $d_\partial$ be the distance to the nearer endpoint. Near that
endpoint put, with the constant $d_0$ of
Lemma~\ref{lem:tail-flat},
\begin{equation}\label{eq:readable-physical-weights}
 x_\partial=\sqrt{\log\frac{d_0}{d_\partial}},\quad
 H=e^{-2x_\partial^3/3},\quad
 E_0=e^{1+x_\partial^{5/2}\log(e+x_\partial)},\quad
 M=e^{1+x_\partial^{5/2}\log^2(e+x_\partial)},
\end{equation}
and extend these functions smoothly and positively through the middle
of the annulus. By \eqref{eq:geometry-endpoint-distance},
$x_\partial=n+O(1)$ on cell $n$. Hence $H$ agrees with $H_n$, and
$E_0$ with $\mathcal E_n$, up to coefficient powers, and $M$ is the
weight \eqref{eq:M-weight}. For every fixed $A$ and $\epsilon>0$,
\begin{equation}\label{eq:full-material-envelope}
 E_0^A\le C_A\,M,\qquad H^\epsilon M^A\le C_{A,\epsilon}.
\end{equation}
The extra logarithm in $M$ is deliberate. However many coefficients
appear after differentiating a fixed expression, their \emph{whole
product} is bounded by a power of $E_0$ and thus costs one power of
$M$, a cost that does not grow with the number of derivatives. And
the cubic exponent of $H$ still dominates every power of $M$.

\paragraph{A unit-scale derivative.}
Near the ends put
\[
 \ell=\frac{d_\partial}{1+x_\partial},\qquad
 d\varrho=\frac{dr}{\ell},\qquad \Dr=\partial_\varrho=\ell\,\partial_r ,
\]
with a positive extension in between. The variable $\varrho$ ranges
over $\R$, and
\[
 |\Dr x_\partial|=\frac1{2x_\partial(1+x_\partial)},\qquad
 |\Dr\log H|=\frac{x_\partial}{1+x_\partial}.
\]
All $\Dr$-derivatives of positive order of $\log H$, $\log M$ and
$\log\ell$ are bounded: in the variable $\varrho$ the weights vary on
the unit scale. An ordinary radial derivative $\partial_r$ costs a
factor bounded by a power of $E_0$.

\paragraph{The exponent $\beta$.}
We fix $\beta=\tfrac34$. What matters is
\begin{equation}\label{eq:beta-choice}
 \tfrac12<\beta<1 .
\end{equation}
The upper inequality leaves room for two neighbouring exponents
$\beta\pm\epsilon$ in $(0,1)$, at which the reference inverse is
available. The lower inequality makes a product of two quantities of
size $H^\beta$ smaller than $H$; this gives the improved estimate for
the metric in Theorem~\ref{thm:new-annular-height}.

\begin{definition}[Spaces of heights, sources and circle data]
\label{def:readable-physical-grading}
All functions are real, even in $\theta$ and $\pi$-periodic. For
$\gamma>0$ and integers $s\ge0$ put
\begin{equation}\label{eq:full-physical-norm}
 \|u\|_{\gamma,s}^2
 =\sum_{a+b+c\le s}\bigl\|M^c\,\Dr^a\partial_\theta^b
                  (H^{-\gamma}u)\bigr\|_{L^2(d\varrho\,d\theta)}^2 .
\end{equation}
\begin{itemize}
\item $\mathcal X^0_\gamma$ is the space of functions with all these
 norms finite.
\item $\mathcal X_\gamma=\mathcal X^0_\gamma\oplus\R\chi_L$, where
 $\chi_L$ is a fixed function of $r$, equal to one near the left
 endpoint and to zero near the right one, whose derivative is
 supported in a compact subset of a positive component. The
 coefficient of $\chi_L$ is the difference of the two endpoint values
 of a height.
\item $\mathcal Y_\beta=\{f\in\mathcal X^0_\beta:\Qint f=0\}$ is the
 space of compatible sources, where
 $\Qint f=\int_{r_-}^{r_+}\av f\,dr$.
\item For a sequence $a=(a_m)$ of functions on the circle, one for
 each negative component,
 \begin{equation}\label{eq:weak-circle-spaces}
  \|a\|_{\mathcal Z_{\gamma,s}}^2
  =\sum_m\sum_{b+c\le s}M_m^{2c}H_m^{-2\gamma}
                    \|\partial_\theta^b a_m\|_{L^2(\T)}^2 ,
 \end{equation}
 and $\mathcal Z_\gamma$ is the space with all these norms finite.
 Its subspace of sequences of constants is $\mathcal K_\gamma$, and
 $\mathcal Z_{\gamma,\mathrm{odd}}$, $\mathcal Z_{\gamma,\mathrm{even}}$
 are its subspaces of odd and of even functions.
\end{itemize}
When no exponent is displayed, a norm of a height or of a source is
the one with $\gamma=\beta$.
\end{definition}

Note that the index $s$ counts derivatives and powers of $M$
together. All these spaces consist of smooth functions which are flat
at the endpoints; the exponent $\gamma$ only measures how fast they
decay.

\paragraph{Tame spaces and tame maps.}
We use the terminology of \cite[Part~II]{Hamilton1982}. A
\emph{graded} Fr\'echet space is one whose topology is defined by an
increasing sequence of norms $\|\cdot\|_0\le\|\cdot\|_1\le\cdots$. A
linear map $L$ between graded spaces is \emph{tame} if there is a
number $d$, its \emph{loss}, with $\|Lf\|_s\le C_s\|f\|_{s+d}$ for
all $s$. For a Banach space $B$, let $\Sigma(B)$ be the space of
sequences $(f_k)_{k\ge0}$ in $B$ for which all the norms
$\|f\|_s=\sum_ke^{sk}\|f_k\|_B$ are finite. A graded space $F$ is
\emph{tame} if there are tame linear maps $J:F\to\Sigma(B)$ and
$P:\Sigma(B)\to F$ with $PJ=I$. A map $\mathcal P$ on an open subset
of a graded space is \emph{tame} if it is continuous and every point
has a neighbourhood on which
$\|\mathcal P(f)\|_s\le C_s(1+\|f\|_{s+d})$ for all $s$; it is
\emph{smooth tame} if it is smooth and all its derivatives are tame.
Since a derivative of order $k$ is linear in its $k$ increments, this
amounts to estimates of the form
\begin{equation}\label{eq:tame-derivative-shape}
 \bigl\|D^k\mathcal P(f)[h_1,\ldots,h_k]\bigr\|_s
 \le C_{s,k}\Bigl[\,\sum_{i}\|h_i\|_{s+d_k}\prod_{j\ne i}\|h_j\|_{b}
       +(1+\|f\|_{s+d_k})\prod_{j}\|h_j\|_{b}\Bigr]
\end{equation}
with a fixed low order $b$. \emph{In every term only one factor
carries the high order.} All estimates of this part have this shape.
Two features will be checked each time: the loss does not depend on
$s$, and the neighbourhood is defined by norms of one fixed order.

The following elementary lemma produces estimates of this shape. It
replaces the usual interpolation inequalities for sup norms by
inequalities for the Fourier transform, which can be verified in a
few lines.

\begin{lemma}[Estimates with one high factor]\label{lem:one-high-factor}
Let $\mathsf C$ be the circle $\T$ or the cylinder $\R\times\T$, let
$\|\cdot\|_{H^s}$ denote the Sobolev norms on $\mathsf C$, and let
$|v|_{\mathcal F}=\|\widehat v\|_{L^1}$ be the $L^1$ norm of the
Fourier transform. Then
$\sup|v|\le|v|_{\mathcal F}\le C\,\|v\|_{H^b}$, where $b=1$ on the
circle and $b=2$ on the cylinder.
\begin{enumerate}
\item \textup{(Interpolation.)} If $k\le i\le m$ and the point $(i,j)$
 lies on the segment from $(k,l)$ to $(m,n)$, then
 \[
  \|f\|_{H^i}\,\|g\|_{H^j}
  \le\|f\|_{H^k}\,\|g\|_{H^l}+\|f\|_{H^m}\,\|g\|_{H^n}.
 \]
\item \textup{(Products.)} If $\alpha_1,\ldots,\alpha_k$ are
 multi-indices with $|\alpha_1|+\cdots+|\alpha_k|\le s$, then
 \[
  \Bigl\|\prod_{i=1}^k\partial^{\alpha_i}v_i\Bigr\|_{L^2}
  \le C_k\sum_{i=1}^k\|v_i\|_{H^s}\prod_{j\ne i}|v_j|_{\mathcal F}.
 \]
 In particular
 $\|v_1v_2\|_{H^s}\le C_s\bigl(\|v_1\|_{H^s}|v_2|_{\mathcal F}
 +|v_1|_{\mathcal F}\|v_2\|_{H^s}\bigr)$.
\item \textup{(Functions.)} Let $\Phi$ be smooth on a neighbourhood of
 the convex hull of $\{0\}$ and the closure of the range of a real
 function $v$. Then
 $\|\Phi(v)-\Phi(0)\|_{H^s}\le C_s\,\|v\|_{H^s}$ for every $s$,
 where $C_s$ depends only on $|v|_{\mathcal F}$ and on the suprema of
 $|\Phi'|,\ldots,|\Phi^{(\max\{1,s\})}|$ on that convex hull.
 The constants are uniform when these derivative bounds and an upper bound
 for $|v|_{\mathcal F}$ are fixed. The same holds for a smooth function
 of several real functions, using a convex region containing zero and the
 closed joint range, with bounds on its derivatives of orders $1$ through
 $\max\{1,s\}$ on that region.
\item \textup{(Composition and inversion on the circle.)} Let $h$ be
 the lift of a diffeomorphism of the circle with $h'\ge c_0>0$. Then,
 for $s\ge1$,
 \[
  \|u\circ h\|_{H^s}
  \le C_s\bigl(\|u\|_{H^s}+\|u\|_{H^3}\,\|h'-1\|_{H^{s-1}}\bigr),
  \qquad
  \|(h^{-1})'-1\|_{H^{s-1}}\le C_s\,\|h'-1\|_{H^{s-1}},
 \]
 where $C_s$ depends only on $c_0$ and $\|h'-1\|_{H^3}$.
\end{enumerate}
\end{lemma}

\begin{proof}
The first assertion is the Cauchy--Schwarz inequality for
$\widehat v=\langle\zeta\rangle^{-b}\cdot\langle\zeta\rangle^b\widehat v$,
since $\langle\zeta\rangle^{-2b}$ is summable over the frequencies of
$\mathsf C$ exactly when $2b$ exceeds its dimension.

(i) Let $i=\vartheta k+(1-\vartheta)m$; then
$j=\vartheta l+(1-\vartheta)n$. By H\"older's inequality for the
Fourier transform,
$\|f\|_{H^i}\le\|f\|_{H^k}^{\vartheta}\|f\|_{H^m}^{1-\vartheta}$, and
likewise for $g$. The product of the two right sides is
$X^\vartheta Y^{1-\vartheta}\le X+Y$ with
$X=\|f\|_{H^k}\|g\|_{H^l}$ and $Y=\|f\|_{H^m}\|g\|_{H^n}$.

(ii) The Fourier transform of the product at the frequency $\zeta$ is
bounded by the integral of
$\prod_i|\zeta_i|^{|\alpha_i|}|\widehat v_i(\zeta_i)|$ over
$\zeta_1+\cdots+\zeta_k=\zeta$. Since
$\langle\zeta_i\rangle\ge1$ and the exponents add up to at most $s$,
\[
 \prod_i|\zeta_i|^{|\alpha_i|}
 \le\max_i\langle\zeta_i\rangle^{s}
 \le\sum_i\langle\zeta_i\rangle^{s}.
\]
So the Fourier transform of the product is bounded by a sum of $k$
convolutions, in each of which one factor is
$\langle\zeta\rangle^s|\widehat v_i|\in L^2$ and the others are
$|\widehat v_j|\in L^1$. Young's inequality and Plancherel's theorem
give the claim.

(iii) A derivative $\partial^\alpha$ of order between $1$ and $s$ of
$\Phi(v)$ is a sum of terms
$\Phi^{(m)}(v)\prod_{i=1}^m\partial^{\alpha_i}v$ with
$\alpha_1+\cdots+\alpha_m=\alpha$; bound $\Phi^{(m)}(v)$ by its
supremum and apply (ii). In $L^2$,
$|\Phi(v)-\Phi(0)|\le\sup|\Phi'|\,|v|$, where this supremum controls
the entire segment from zero to each input value. This first-derivative
bound is also required when $s=0$.

(iv) Change of variables gives
$\|f\circ h\|_{L^2}\le c_0^{-1/2}\|f\|_{L^2}$. For $s\le3$ the chain
rule expresses the derivatives of $u\circ h$ through
$u^{(m)}\circ h$, $m\le s$, and bounded functions, which proves the
first estimate without its second term. Let $s\ge4$ and assume the
estimate for $s-1$. Since $(u\circ h)'=(u'\circ h)\,h'$, part~(ii) on
the circle gives
\[
 \|u\circ h\|_{H^s}
 \le C\bigl(\|u\|_{L^2}+\|u'\circ h\|_{H^{s-1}}\,\|h'\|_{H^1}
 +\|u'\circ h\|_{H^1}\,(1+\|h'-1\|_{H^{s-1}})\bigr).
\]
By the induction hypothesis
$\|u'\circ h\|_{H^{s-1}}\le C(\|u\|_{H^s}+\|u\|_{H^4}\|h'-1\|_{H^{s-2}})$,
and by (i), applied to the points $(3,s-1)$, $(4,s-2)$ and $(s,2)$,
\[
 \|u\|_{H^4}\,\|h'-1\|_{H^{s-2}}
 \le\|u\|_{H^3}\,\|h'-1\|_{H^{s-1}}+\|u\|_{H^s}\,\|h'-1\|_{H^{2}} .
\]
This proves the first estimate. For the inverse $\iota=h^{-1}$ we
have $\iota'-1=G\circ\iota$ with $G=-(h'-1)/h'$, a smooth function of
$h'-1$. By the chain rule $\iota'$ is bounded below by $1/\sup h'$
and $\|\iota'-1\|_{H^3}$ is bounded in terms of $c_0$ and
$\|h'-1\|_{H^3}$; and
$\|G\|_{H^k}\le C\|h'-1\|_{H^k}$ by (iii). The first estimate, with
$\iota$ in the place of $h$, gives
\[
 \|\iota'-1\|_{H^{s-1}}
 \le C\bigl(\|G\|_{H^{s-1}}+\|G\|_{H^3}\,\|\iota'-1\|_{H^{s-2}}\bigr),
\]
and induction on $s$ proves the second estimate.
\end{proof}

Functions on a cylinder with boundary, such as the pieces of
Section~\ref{sec:stable-inverse}, are first extended to
$\R\times\T$ by a fixed extension operator which is bounded in all
Sobolev norms; parts (i)--(iii) then apply to them as well.

\begin{lemma}[Flows on the circle]\label{lem:circle-flow}
Let $G(t,\theta)$ be smooth and periodic in $\theta$, and let
$\Theta$ solve $\partial_t\Theta=G(t,\Theta)$ for $0\le t\le1$ with
$\Theta(0,\theta)=\theta$. Put
$\Lambda_s=\int_0^1\|G_\theta(t,\cdot)\|_{H^s}\,dt$. Then
\[
 \Theta_\theta=\exp\int_0^tG_\theta(\tau,\Theta(\tau,\theta))\,d\tau
 \in\bigl[e^{-C\Lambda_1},e^{C\Lambda_1}\bigr],
 \qquad
 \sup_t\|\Theta_\theta(t,\cdot)-1\|_{H^s}\le C_s\,\Lambda_s ,
\]
where $C_s$ depends only on $\Lambda_4$. If $G_*$ does not depend on
$\theta$ and $\Theta_*=\theta+\int_0^tG_*$ is its flow, then
$\sup_t\|\Theta-\Theta_*\|_{H^s}
\le C_s\int_0^1\|(G-G_*)(t,\cdot)\|_{H^s}\,dt$, with $C_s$ depending
only on $\Lambda_4$ and on $\int_0^1\|(G-G_*)(t,\cdot)\|_{H^3}\,dt$.
\end{lemma}

\begin{proof}
The formula for $\Theta_\theta$ follows from
$\partial_t\Theta_\theta=(G_\theta\circ\Theta)\,\Theta_\theta$, and
the bounds from $|G_\theta|\le C\|G_\theta\|_{H^1}$. The derivatives
$\partial_\theta^k\Theta$, $2\le k\le4$, solve linear equations with
the same integrating factor, whose sources are polynomials in the
lower derivatives and in $(\partial_\theta^mG)\circ\Theta$,
$2\le m\le4$; hence they are bounded in terms of $\Lambda_4$. Now let
$y=\Theta_\theta-1$, so that
$\partial_ty=(G_\theta\circ\Theta)(1+y)$. By
Lemma~\ref{lem:one-high-factor}(ii) and~(iv), with the low norms just
bounded,
\[
 \|\partial_ty\|_{H^s}
 \le C_s\bigl(\|G_\theta(t,\cdot)\|_{H^s}
 +\|G_\theta(t,\cdot)\|_{H^3}\,\|y\|_{H^s}\bigr),
\]
and Gronwall's inequality gives the second assertion. For the last
one, $\partial_t(\Theta-\Theta_*)=(G-G_*)(t,\Theta)$; apply
Lemma~\ref{lem:one-high-factor}(iv) and integrate.
\end{proof}

The point of the last statement is that a flow with a large speed
which does not depend on the angle, as for the reference profile, does
not enter the constants: only the angular derivative of the speed
does.

\begin{lemma}[Calculus and flat extension]
\label{prop:full-physical-tame}
The spaces of Definition~\ref{def:readable-physical-grading}, graded by
$s$, are tame Fr\'echet spaces in the sense of
\cite{Hamilton1982}. Products, differential expressions of fixed
order, and reciprocals of coefficients, after their reference factors
have been separated, are smooth tame maps on neighbourhoods on which
the remaining denominators are bounded below; the loss of derivatives
does not depend on the order $s$ of the output. Every element of
$\mathcal X^0_\gamma$, $\gamma>0$, extends smoothly by zero across
both endpoints, and so do its derivatives with respect to parameters.
\end{lemma}

\begin{proof}
\step{1: tameness}
Write $v=H^{-\gamma}u$. Choose functions $\phi_i(\varrho)$, $i\in\Z$,
with bounded derivatives, finite overlap and $\sum_i\phi_i^2=1$,
supported on intervals of unit length. Since $\log M$ has bounded
derivatives, $M\asymp M_i:=M(i)$ on the support of $\phi_i$, and the
norm of $v$ is equivalent to
\[
 \sum_{i,j}\int(M_i^2+\xi^2+j^2)^s\,
                    |\widehat{\phi_iv}(\xi,j)|^2\,d\xi .
\]
Let $\omega_i=(M_i^2+\xi^2+j^2)^{1/2}$. A smooth partition
$\zeta_k$ in the variable $\log\omega_i$ gives maps $J$ and $P$ with
$PJ=I$ to and from the space of exponentially decreasing sequences in
$\ell^2_iL^2_\xi\ell^2_j$: namely
$(Jv)_{k,i}=\zeta_k(\omega_i)\,\widehat{\phi_iv}$ and
$Pb=\sum_i\phi_i\,\mathcal F^{-1}\sum_k\widetilde\zeta_k(\omega_i)\,b_{k,i}$,
where $\widetilde\zeta_k$ equals one on the support of $\zeta_k$. By
the Cauchy--Schwarz inequality in $k$,
$\|Jv\|_s\le C_s\|v\|_{s+1}$, and by finite overlap
$\|Pb\|_s\le C_s\|b\|_s$. This exhibits the space as a tame direct
summand of a model space, and provides smoothing operators. Parity
projections preserve the construction, and the scalar multiple of
$\chi_L$ is a finite-dimensional summand. If $\psi$ is compactly
supported with $\Qint\psi=1$, then $f\mapsto f-\psi\,\Qint f$ is a
projection onto $\mathcal Y_\beta$; the functional $\Qint$ is
continuous because it is the pairing of $H^{-\beta}f$ with the square
integrable function $H^\beta\ell$. The same Fourier construction
applies to the spaces of circle data.

\step{2: products}
Let $v_k=H^{-\gamma_k}u_k$, $k=1,2$, so that
$H^{-\gamma_1-\gamma_2}u_1u_2=v_1v_2$. On the support of $\phi_i$ the
norm of order $s$ of $v$ is equivalent to
$\sum_{c\le s}M_i^c\,\|\phi_iv\|_{H^{s-c}}$. Write
$\phi_iv_1v_2=(\phi_iv_1)(\widetilde\phi_iv_2)$ with a cutoff
$\widetilde\phi_i$ equal to one on the support of $\phi_i$, apply
Lemma~\ref{lem:one-high-factor}(ii) on the cylinder, and attach the
power $M_i^c$ to the factor of high order. Summing the squares over
$i$, with the low factor estimated uniformly in $i$, gives
\begin{equation}\label{eq:weighted-product}
 \|u_1u_2\|_{\gamma_1+\gamma_2,s}
 \le C_s\bigl(\|u_1\|_{\gamma_1,s}\,\|u_2\|_{\gamma_2,2}
 +\|u_1\|_{\gamma_1,2}\,\|u_2\|_{\gamma_2,s}\bigr) .
\end{equation}
Weights add under multiplication, and only one factor carries the
high order. Products of several factors, and products of derivatives
of total order at most $s$, are estimated in the same way.

\step{3: coefficients}
Every coefficient that occurs below is a product of a \emph{reference
function}, which depends only on $r$ and whose derivatives are
bounded by coefficient powers, and a function
$\Phi(\varepsilon_1,\ldots,\varepsilon_p)$, where $\Phi$ is smooth
near the origin and each $\varepsilon_k$ is a reference function
times a derivative of positive order of a height, a control term, or
a displacement of a chart. Examples are
\[
 \frac w{w_*}=e^{\nu\chi_*+\sum\kappa_m\chi_m},\qquad
 \frac Qw=1+\frac{Z_r+EZ_{\theta\theta}/r}{w},\qquad
 D=1-2wZ_r-Z_r^2-E\Bigl(\frac{Z_\theta}r\Bigr)^2,
\]
and the reciprocals and powers of these quantities that enter
\eqref{eq:linear-coefficients}. On a neighbourhood of low order the
$\varepsilon_k$ are small in the sup norm, so $\Phi$ is evaluated
where it is smooth. By Lemma~\ref{lem:one-high-factor}(iii), applied
on the supports of the $\phi_i$, the norm of order $s$ of
$\Phi(\varepsilon)-\Phi(0)$ is bounded by the norms of order $s$ of
the $\varepsilon_k$, with a constant that depends on their low norms
only. Moreover
$\Phi(\varepsilon)-\Phi(0)=\sum_k\varepsilon_k\Phi_k(\varepsilon)$
with smooth $\Phi_k$, so this difference carries the weight
$H^\gamma$ of the $\varepsilon_k$. The same holds for the derivatives
with respect to the state, which are of the same form. In
particular, reciprocals of the denominators $Q$, $D$ and $1-p^2$ are
smooth tame after the reference factor $w_*$ has been separated.

After a complete expression has been expanded, all its reference
functions together are bounded by $E_0^{C_s}$, which costs one order
by \eqref{eq:full-material-envelope}. The expressions have fixed
differential order, so the total loss is fixed. Restriction to a
circle $r=\text{const}$ costs one radial derivative, and a fixed
cutoff provides a right inverse.

\step{4: flat extension}
The operator $\partial_r^a$ is a combination of $\Dr$-derivatives of
order at most $a$ with coefficients bounded by $E_0^{C_a}$. By the
Sobolev embedding on unit intervals,
$d_\partial^{-k}\,|\partial_r^a\partial_\theta^bu|$ is bounded by a
seminorm of $u$ times $H^\gamma E_0^{C_{a,b,k}}$, which tends to zero
at the endpoints. The same holds after differentiation with respect to
parameters.
\end{proof}

\subsection{The geometry at the roots, and the travel-angle controls}
\label{subsec:root-geometry}

\paragraph{Scales in a cell.}
Recall from Section~\ref{sec:radial-tails} that on cell $n$
\[
 A_n\asymp e^{-n^2},\quad B_n\asymp A_n^9,\quad
 c\,\mathfrak r_nA_n^5\le r_X\le C\,\mathfrak r_nA_n,\quad
 w_*=\lambda_n\,r_X,\quad \log\lambda_n=-n^{5/2}+O(n),
\]
with $\mathfrak r_n$ bounded above and below, that
$r_X=\partial r/\partial X$ satisfies $r_{XX}=-\rho\,r$ with
$\rho=A_n(e^{-X}-e^{-2X})-B_n$, that $cA_{n+1}\le d_\partial\le CA_n$,
and that every fixed derivative of $w_*^{\pm1}$ costs a coefficient
power (Lemma~\ref{lem:tail-flat}). These comparisons persist on a
bounded enlargement of a cell.

\paragraph{The characteristic speed near a root.}
Let $x$ be the distance from a root of a negative component, counted
positively into the component and negatively on the positive side. On
a cell, $w_*=\lambda_nr_X$ and $r_{XX}=-\rho\,r$ give
$w_*'=-\lambda_n\rho\,r/r_X$, so the reference characteristic speed
\eqref{eq:completed-characteristics} is
\[
 v_0=\sqrt{-\frac{w_*'}{rw_*}}=\frac{\sqrt\rho}{r_X}
 =\sqrt x\;c_i(x),\qquad i\in\{L,R\},
\]
where $c_i$ is smooth and positive on a neighbourhood of the root,
on both sides.

\begin{lemma}[Root collars]\label{lem:quantitative-root-collars}
At the left and right roots of the negative component of cell $n$,
\begin{equation}\label{eq:geometry-root-scales}
 c_L(0)\asymp\mathfrak r_n^{-3/2}A_n^{-1},\qquad
 c_R(0)\asymp\mathfrak r_n^{-3/2}A_n^{-3},
\end{equation}
and for every $k$, on the interval $|x|\le c\,\mathfrak r_nA_n^5$,
\begin{equation}\label{eq:geometry-normalized-root-jets}
 \Bigl|\partial_r^k\Bigl(\frac{c_i}{c_i(0)}\Bigr)^{\pm1}\Bigr|
 \le C_k\,(\mathfrak r_nA_n^5)^{-k}.
\end{equation}
Fix a sufficiently small $\epsilon_{\rm root}>0$ and put
\begin{equation}\label{eq:geometry-root-collar-width}
 \delta_n^{\rm root}=\epsilon_{\rm root}\,\mathfrak r_n\,A_n^{20}.
\end{equation}
Then the collars $|x|<8\delta_n^{\rm root}$ of the two roots are
disjoint and lie inside the cell, and on them
$c_i(x)/c_i(0)\in[\tfrac12,2]$. The finitely many roots of the central
model have collars with the same properties and fixed constants.
\end{lemma}

\begin{proof}
\step{1: the value at the root}
Let $X_i$ be the value of the model coordinate at the root. Since
$v_0^2=\rho/r_X^2$ and $x=|r-r(X_i)|$, division by $x$ and passage to
the limit give
\[
 c_i(0)^2=\frac{|\rho_X(X_i)|}{r_X(X_i)^3}.
\]
With $t=e^{-X}$ the roots solve $A_n(t-t^2)=B_n$; hence
$e^{-X_L}=1+O(A_n^8)$ and $e^{-X_R}\asymp A_n^8$, and
$|\rho_X(X_L)|\asymp A_n$, $|\rho_X(X_R)|\asymp A_n^9$. By
Lemma~\ref{lem:cell} and the choice of the solution in
Section~\ref{subsec:one-cell}, $r_X(X_L)\asymp\mathfrak r_nA_n$ and
$r_X(X_R)=\epsilon_R\,\mathfrak r_nA_n^5$. This gives
\eqref{eq:geometry-root-scales}.

\step{2: derivatives}
Division of a function with a simple zero by the distance to the zero
is the operation $f(s)/s=\int_0^1f'(ts)\,dt$; it costs one derivative
and preserves bounds. On $|X-X_i|\le c$ the quotient
$\rho/(X_i-X)$, with the sign which makes it positive, is comparable
to $|\rho_X(X_i)|$, and its derivatives in $X$ are bounded by
multiples of it; this is read off from the explicit function $\rho$.
At a root $r_{XX}=0$ and $r_{XXX}=-\rho_Xr$, and the differential
equation shows that $r_X$ is comparable to $r_X(X_i)$ on
$|X-X_i|\le c$, with derivatives in $X$ bounded by multiples of it.
The same then holds for the quotient $(r-r(X_i))/(X-X_i)$. Since
\[
 c_i^2=\frac{\rho}{X_i-X}\cdot\frac{X_i-X}{r(X_i)-r}\cdot\frac1{r_X^2},
\]
with the signs chosen as above, and
$\partial_r=r_X^{-1}\partial_X$ with $r_X\ge c\,\mathfrak r_nA_n^5$,
we obtain \eqref{eq:geometry-normalized-root-jets} on the image of
$|X-X_i|\le c$, an interval of radial length at least
$c\,\mathfrak r_nA_n^5$ on either side of the root. At the left root
the same argument gives the better scale $\mathfrak r_nA_n$.

\step{3: collars}
The width $8\delta_n^{\rm root}$ divided by $\mathfrak r_nA_n^5$ is
$8\epsilon_{\rm root}A_n^{15}$. So for small $\epsilon_{\rm root}$ the
case $k=1$ of \eqref{eq:geometry-normalized-root-jets} keeps
$c_i(x)/c_i(0)$ in $[\tfrac12,2]$ on the collar. The collars lie in
the cell and are disjoint because each part of the cell has radial
length comparable to $A_n$. The central model has finitely many
roots, at which $c_i(0)>0$, and fixed collars are chosen there. Note
that the collar was chosen before any order of differentiation; its
width, its reciprocal, and the derivatives of cutoffs adapted to it
are coefficient powers.
\end{proof}

\begin{lemma}[Travel-angle controls]
\label{lem:quantitative-optical-controls}
For a profile $w$ let
$T_n(w)=\int\sqrt{-w'/(rw)}\,dr$ be the travel angle of its $n$th
negative component. There are functions $\chi_n$, supported in
$\delta_n^{\rm root}<x<2\delta_n^{\rm root}$ near the left root of
that component, such that
\[
 \frac{d}{dt}\Big|_{t=0}T_n\bigl(w_*e^{t\chi_n}\bigr)=1,
 \qquad \|\partial_r^k\chi_n\|_\infty
       \le C_k\,\mathfrak r_n^{-k}A_n^{9-20k}.
\]
There is also a cutoff function of $r$ which rises from $0$ to $1$ on
$3\delta_n^{\rm root}<x<4\delta_n^{\rm root}$, with the corresponding
bounds. On the supports of $\chi_n$ and of the derivative of this
cutoff,
\[
 v_0\asymp\mathfrak r_n^{-1}A_n^9,\qquad
 \Bigl|\frac{w_*'}{w_*}\Bigr|\asymp\mathfrak r_n^{-1}A_n^{18}.
\]
\end{lemma}

\begin{proof}
Let $\phi\ge0$ be a nonzero smooth function supported in $(1,2)$ and
put $f_n=\phi(x/\delta_n^{\rm root})$. For the variation
$w_*e^{tf_n}$ the speed changes by $\delta v_0=-f_n'/(2rv_0)$, so
\[
 \frac{d}{dt}\Big|_{t=0}T_n
 =\frac12\int\Bigl(\frac1{rv_0}\Bigr)'f_n\,dr
 =-\frac{(\delta_n^{\rm root})^{-1/2}}{4\,r_{n,L}\,c_L(0)}
       \int_1^2t^{-3/2}\phi(t)\,dt\;\bigl(1+O(A_n^{19})\bigr),
\]
where $r_{n,L}$ is the radius of the root. This is nonzero and
comparable to $-A_n^{-9}$. Let $\chi_n$ be $f_n$ divided by this
number; the bounds follow by scaling. The bounds for $v_0$ follow
from \eqref{eq:geometry-root-scales}, and $|w_*'/w_*|=rv_0^2$. The
cutoff is the integral of a rescaled bump.
\end{proof}

\subsection{The state and one neighbourhood for all components}
\label{subsec:neighbourhood}

The state space is
\begin{equation}\label{eq:rewrite-state}
 \bx=(u,\nu,\kappa,h_L),\qquad
 u\in\Phi_*+\ker\ell_*\subset\mathcal X_\beta,\quad
 \nu\in\R,\quad
 \kappa\in\mathcal K_\beta,\quad
 h_L-\id\in\mathcal Z_{\beta,\mathrm{odd}} .
\end{equation}
Each $h_{L,m}$ is an increasing map of $\R$ which commutes with the
translation by $\pi$, that is, the lift of a diffeomorphism of the
circle $\R/\pi\Z$. The normalization $\ell_*u=1$ removes the freedom
of adding a multiple of the kernel $\Phi_*$. The size of a state,
together with the parameter, is measured by the seminorms
\begin{equation}\label{def:joint-displacement-seminorms}
 \mathfrak n_s=|q|+\|u-\Phi_*\|_s+|\nu|
           +\|\kappa\|_{\mathcal K_{\beta,s}}
           +\|h_L-\id\|_{\mathcal Z_{\beta,s}} .
\end{equation}
The unscaled $u$ appears here, and not only $qu$, because the
derivative of $\mathcal F_0(w,u)=\mathscr L_wu$ with respect to the
metric depends on $u$ even at $q=0$.

We shall also need localized versions. For a component $m$, positive
or negative, consider the component together with the neighbouring
components that enter the construction of its coordinates; for a
positive component these are the two adjacent negative ones. These
neighbourhoods have bounded overlap. Let $\mathfrak g_{s,m}$ be the
norm of order $s$ of $(qu,\nu,\kappa,h_L-\id)$ restricted to this
neighbourhood, where only derivatives of positive order of the height
are used, and compact scalars are counted only where their supports
meet the neighbourhood. Then
$(\sum_m\mathfrak g_{s,m}^2)^{1/2}\le C_s\,\mathfrak n_s$.

For the joint problem we also use the localized seminorm
$\mathfrak n_{s,m}$ of
$(q\,\Phi_*,\,u-\Phi_*,\,\nu,\,\kappa,\,h_L-\id)$, defined in the
same way. For $|q|\le1$ the identity $qu=q\Phi_*+q(u-\Phi_*)$ gives
$\mathfrak g_{s,m}\le C_s\,\mathfrak n_{s,m}$, and
$(\sum_m\mathfrak n_{s,m}^2)^{1/2}\le C_s\,\mathfrak n_s$.

\begin{lemma}[One neighbourhood for the type and the denominators]
\label{lem:common-geometric-neighborhood}
There are a finite order $b_0$ and an $\epsilon_0>0$ such that on the
neighbourhood $\mathfrak n_{b_0}<\epsilon_0$ the profile $w$ has the
same positive and negative components and the same germs at the roots
as $w_*$, and, with $Z=qu$ and $E=1+w^2$,
\[
 \frac{w}{w_*}\in\bigl[\tfrac12,2\bigr],\qquad
 \frac Qw\in\bigl[\tfrac78,\tfrac98\bigr],\qquad
 D\in\bigl[\tfrac12,\tfrac32\bigr],\qquad
 1-\Bigl(\frac{Z_\theta}{r}\Bigr)^2\ge\tfrac34
\]
on the whole annulus. Here
\begin{equation}\label{eq:geometry-low-QD}
 Q=w+Z_r+\frac{E\,Z_{\theta\theta}}r,\qquad
 D=1-2wZ_r-Z_r^2-E\Bigl(\frac{Z_\theta}r\Bigr)^2
\end{equation}
are the quantities \eqref{eq:height-jet-notation} for $z=h+Z$. The
neighbourhood does not depend on any order of differentiation.
\end{lemma}

\begin{proof}
Only finitely many derivatives of low order enter these conditions.

\step{1: low-order bounds on all components}
On the neighbourhood, $\|u\|_{b_0}\le\|\Phi_*\|_{b_0}+\epsilon_0$ and
$|q|<\epsilon_0$, so the geometric displacement
$(Z,\nu,\kappa,h_L-\id)$ is of size $C\epsilon_0$ in the norms of
order $b_0$. Let $K$ be the largest order of a derivative that occurs
in the conditions, and take $b_0\ge K+2$. The Sobolev embedding on
unit intervals of the variable $\varrho$, and the conversion of
$\Dr$ into $\partial_r$ at the cost of coefficient powers, give on
component $m$
\begin{equation}\label{eq:geometry-low-jets}
 \max_{1\le a+b\le K}|\partial_r^a\partial_\theta^bZ|
 +|\kappa_m|\,\|\chi_m\|_{C^K}+\|h_{L,m}-\id\|_{C^K}
 \le C\,\epsilon_0\,H_m^\beta\,\mathcal E_m^{C}.
\end{equation}
Only derivatives of positive order of $Z$ occur: the endpoint constant
of a height does not decay, but it does not enter any coefficient.
The compact control $\nu\chi_*$ obeys the same bound on the finitely
many components it meets. Now
\begin{equation}\label{eq:geometry-smallness-absorption}
 \sup_mH_m^\beta\,\mathcal E_m^{A}<\infty
 \qquad\text{and}\qquad
 H_m^\beta\,\mathcal E_m^{A}\longrightarrow0
 \qquad\text{for every fixed }A,
\end{equation}
because $\log\mathcal E_m=o(S_m)$. Hence one $\epsilon_0$ makes the
right side of \eqref{eq:geometry-low-jets}, multiplied by any fixed
coefficient power, as small as we wish on all components at once.

\step{2: the type}
Write $w=w_*e^f$ with $f=\nu\chi_*+\sum_m\kappa_m\chi_m$. The controls
vanish near every root, so the germs of $w$ at the roots are those of
$w_*$. On the support of $\chi_n$,
Lemma~\ref{lem:quantitative-optical-controls} gives
\[
 \frac{|\kappa_n\chi_n'|}{|w_*'/w_*|}\le C\,|\kappa_n|\,A_n^{-29},
\]
which is small by Step~1. On the support of $\chi_*$ the function
$w_*'/w_*$ has a positive minimum. So $w'/w=w_*'/w_*+f'$ has the sign
of $w_*'/w_*$ away from the roots, and $|f|\le\log2$.

\step{3: the denominators}
By \eqref{eq:geometry-profile-jets}, $w^{-1}=w_*^{-1}e^{-f}$ is
bounded by a coefficient power. So \eqref{eq:geometry-low-QD},
\eqref{eq:geometry-low-jets} and
\eqref{eq:geometry-smallness-absorption} make $(Q-w)/w$, $D-1$ and
$Z_\theta/r$ uniformly small.
\end{proof}

\begin{proposition}[Uniform geometric calculus]
\label{prop:uniform-physical-geometry}
After $b_0$ is increased and $\epsilon_0$ decreased, the following
hold on the neighbourhood $\mathfrak n_{b_0}<\epsilon_0$.
\begin{enumerate}
\item \textup{(Reference.)} The coordinates, inverse coordinates,
 Jacobians, height factors, source factors and potentials constructed
 in
 Sections~\ref{subsec:return-correction}--\ref{subsec:positive-crossing}
 are defined. At the reference state their derivatives of order $s$
 on component $m$, and those of the reciprocals that occur, are
 bounded by $C_s\mathcal E_m^{C_s}$.
\item \textup{(Differences.)} For a general state the difference
 $\Delta c_m$ of any of these quantities from its reference value
 satisfies
 \begin{equation}\label{eq:geometry-current-jet-bound}
  \|\Delta c_m\|_{C^s}
  \le C_s\,\mathcal E_m^{C_s}\,H_m^\beta\;\mathfrak g_{s+d_*,m}
 \end{equation}
 with a fixed $d_*$. In particular the differentials of the
 coordinates and the height factors remain comparable to their
 reference values.
\item \textup{(Divided coefficients and parameters.)} The same bounds
 hold for the differences of the principal coefficients on positive
 components divided by $y^3(1-y)^3$
 (Section~\ref{subsec:positive-crossing}), and for any fixed number
 of derivatives with respect to the state, with estimates of the
 shape \eqref{eq:tame-derivative-shape} in the seminorms
 $\mathfrak n_{s,m}$.
\end{enumerate}
In all these estimates the complete product of reference functions is
bounded by a power of $E_0$ and is absorbed into one power of the
weight $M$.
\end{proposition}

The objects in the proposition are constructed in the next three
subsections; the proof is given in
Section~\ref{subsec:geometry-proofs}. The mechanism is simple. Every
coordinate is built from a characteristic flow. At the reference
state the speed of the flow does not depend on the angle, so by
Lemma~\ref{lem:circle-flow} only the deformation enters the
estimates for the flow, and the deformation carries the factor
$H_m^\beta$, which by \eqref{eq:geometry-smallness-absorption}
dominates every coefficient power.

\subsection{Fixing the return maps by their own residual}
\label{subsec:return-correction}

\paragraph{Characteristics of $\bar L$.}
Start from the operator $\bar L$ of
Lemma~\ref{lem:type-completion}, with coefficients $(a,b,\bar c)$, on
a negative component. By \eqref{eq:completed-characteristics} its
characteristic slopes are
\[
 c_\pm=d\pm v,\qquad d=\frac ba,\qquad
 v=\frac{\sqrt{b^2-a\bar c}}{a}.
\]
Because the height is even in $\theta$, the reflection
$\refl\theta=-\theta$ exchanges the two families:
$c_+=-c_-\circ\refl$. Let $f$ be the \emph{minus return map}: the map
from the left boundary circle of the component to the right one along
the characteristics of slope $c_-$. Let $T$ be the travel angle of
the reference profile on this component. At the reference state
$f=R_{-T}$.

\paragraph{Charts and the return residual.}
Given the chart $h_L$ of the left circle, define
\begin{equation}\label{eq:full-height-one-chart}
 \widetilde h_R=R_{-T}\circ h_L\circ f^{-1},\qquad
 h_R=\Odd\,\widetilde h_R,\qquad
 \mathfrak e=(\Even\,\widetilde h_R)\circ h_R^{-1},
\end{equation}
where $\Odd$ and $\Even$ denote the odd and even parts of a function
of $\theta$. The derivative of $h_R$ is an average of two positive
derivatives, so $h_R$ is increasing; it is the chart of the right
circle. By construction
$\widetilde h_R=(\id+\mathfrak e)\circ h_R$, hence
\begin{equation}\label{eq:return-identity}
 (\id+\mathfrak e)\circ h_R\circ f\circ h_L^{-1}=R_{-T}.
\end{equation}
So $\mathfrak e=0$ means that the minus return map is the rotation
$R_{-T}$ when read in the charts $h_L$ and $h_R$. Since the charts
are odd and the plus return map is $\refl\circ f\circ\refl$, the plus
return map is then $R_{T}$. Thus the single even equation $\mathfrak e=0$
imposes both return conditions. Its mean concerns the travel angle,
its oscillating part the parametrization of the boundary circles.

\paragraph{The derivative with respect to the chart.}
Vary the chart by $\delta h_L=\psi\circ h_L$ with $\psi$ odd, keeping
$f$ fixed. Then
\begin{equation}\label{eq:full-height-chart-derivative}
 \begin{split}
 D_h\mathfrak e[\psi]&=\mathcal M_T\psi+\mathcal Q_h(\mathfrak e,\psi),\qquad
 \mathcal M_T\psi=\Even\bigl[\psi(\theta+T)\bigr],\\
 \mathcal Q_h(\mathfrak e,\psi)&=
 \Even\Bigl[\mathfrak e\int_0^1
       \psi'(\theta+T+t\mathfrak e)\,dt\Bigr]
 -\mathfrak e'\;\Odd\bigl[\psi(\theta+T+\mathfrak e)\bigr].
 \end{split}
\end{equation}
Indeed, the variation of $\widetilde h_R$, composed with $h_R^{-1}$,
is $\psi(\theta+T+\mathfrak e)$; its odd part is the variation of
$h_R$ in the same coordinates, and the formula follows by
differentiating $\mathfrak e\circ h_R=\Even\,\widetilde h_R$. The
term $\mathcal Q_h$ is bilinear in $(\mathfrak e,\psi)$ in the sense
needed later and vanishes when $\mathfrak e=0$. The leading term is
explicit:
\begin{equation}\label{eq:rewrite-circle-inverse}
 \mathcal M_T\bigl(\sin j\theta\bigr)=\sin(jT)\cos j\theta,
 \qquad\text{so}\qquad
 \mathcal M_T^{-1}\bigl(g_j\cos j\theta\bigr)
       =\frac{g_j}{\sin(jT)}\,\sin j\theta
 \quad(j>0\ \text{even}).
\end{equation}
By Proposition~\ref{prop:rad-fixed-theorem}(ii) the inverse loses one
angular derivative and a coefficient power. It does not exist on
constants: the mean of $\mathfrak e$ cannot be corrected by a chart.
This is the reason for the radial controls $\kappa_m$.

\begin{proposition}[The return correction]
\label{prop:optical-residual-completion}
On the neighbourhood of
Proposition~\ref{prop:uniform-physical-geometry}, the operator
$\bar L$ can be replaced by a symmetric operator in divergence form
\[
 \widehat L=\partial_r(a\,\partial_r+\widehat b\,\partial_\theta)
           +\partial_\theta(\widehat b\,\partial_r
                            +\widehat c\,\partial_\theta)
\]
which has the same coefficient $a$, coincides with $\bar L$ on the
positive components and near the roots, and whose two return maps
$\widehat f_\pm$ on every negative component satisfy
\[
 h_R\circ\widehat f_\pm\circ h_L^{-1}=R_{\pm T}
\]
exactly. The difference has the form
\begin{equation}\label{eq:optical-residual-completion}
 \bar L-\widehat L=\mathscr E_{\bx}(\mathfrak e),
\end{equation}
where $(e,\eta)\mapsto\mathscr E_{\bx}(e)\eta$ is bilinear, smooth
tame in all its arguments, and takes values in $\mathcal Y_\beta$.
\end{proposition}

\begin{proof}
\step{1: deforming the flow}
Let $\Phi_r$ be the flow of the minus characteristics from the left
boundary to the radius $r$, so that $f$ is its value at the right
boundary. The return map we want is
$f_*=h_R^{-1}\circ R_{-T}\circ h_L$. With $\mathsf g=h_R\circ f$,
\eqref{eq:return-identity} gives
\[
 K:=f^{-1}\circ f_*=\mathsf g^{-1}\circ(\id+\mathfrak e)\circ\mathsf g,
 \qquad \Delta:=K-\id .
\]
Let $\chi(r)$ be the cutoff of
Lemma~\ref{lem:quantitative-optical-controls}, rising from $0$ to $1$
inside the component, and put
\[
 J_r=\id+\chi(r)\,\Delta,\qquad \widehat\Phi_r=\Phi_r\circ J_r .
\]
Each $J_r$ is increasing, being a convex combination of increasing
maps, and $\widehat\Phi_r$ runs from the identity to $f\circ K=f_*$.
The curves $r\mapsto\widehat\Phi_r(\theta_0)$ have the slope
\begin{equation}\label{eq:readable-optical-velocity}
 \widehat c_-(r,\theta)=c_-(r,\theta)+
 \Phi_r'(\Phi_r^{-1}\theta)\;\chi'(r)\;
                   \Delta\bigl(J_r^{-1}\Phi_r^{-1}\theta\bigr).
\end{equation}

\step{2: the new coefficients}
Define $\widehat c_+=-\widehat c_-\circ\refl$ and
\[
 \widehat d=\frac{\widehat c_++\widehat c_-}2,\quad
 \widehat v=\frac{\widehat c_+-\widehat c_-}2,\qquad
 \widehat b=a\,\widehat d,\quad
 \widehat c=a\,(\widehat d^{\,2}-\widehat v^{\,2}).
\]
Then $\widehat L$ has the characteristic slopes $\widehat c_\pm$, its
minus return map is $f_*$, and its plus return map is the reflection
of $f_*$; these are the asserted returns. On the support of $\chi'$
the speed $v$ is bounded below by
Lemma~\ref{lem:quantitative-optical-controls}, while the correction
in \eqref{eq:readable-optical-velocity} carries the factor
$H_m^\beta\mathcal E_m^C$. So on one low-order neighbourhood the
correction is less than $v/2$ and $\widehat v>0$. Near the roots and
on positive components nothing is changed.

\step{3: the error is bilinear}
From the definition of $K$,
\[
 \Delta(\theta)=\mathfrak e(\mathsf g\theta)
 \int_0^1(\mathsf g^{-1})'\bigl(\mathsf g\theta
          +t\,\mathfrak e(\mathsf g\theta)\bigr)\,dt .
\]
Keep the state and the points at which the functions in this formula
are evaluated, and replace only the leading factor $\mathfrak e$ by an
independent function $e$. This defines increments $d_1[e]$ and
$v_1[e]$ of $d$ and $v$ which are linear in $e$ and equal the actual
increments $d_1=\widehat d-d$, $v_1=\widehat v-v$ when
$e=\mathfrak e$. The coefficients of $\bar L-\widehat L$ are then
$\delta b[\mathfrak e]$ and $\delta c[\mathfrak e]$, where
\[
 \delta b[e]=-a\,d_1[e],\qquad
 \delta c[e]=-a\,(2d+d_1)\,d_1[e]+a\,(2v+v_1)\,v_1[e] ,
\]
and $\mathscr E_{\bx}(e)\eta$ is the corresponding expression in
divergence form applied to $\eta$. It is bilinear in $(e,\eta)$. It
is supported in the interior of the component, so it is flat at the
endpoints and has integral zero. Each of its terms contains the
factor $e$ and a derivative of positive order of $\eta$, both of
weight $H_m^\beta$; after division by the weight of the source space,
a factor $H_m^\beta$ is left, which absorbs every coefficient power.
Lemmas~\ref{lem:one-high-factor} and~\ref{lem:circle-flow}, with
the bounds of
Lemma~\ref{lem:uniform-characteristic-flow-constants}, give the tame
estimates.
\end{proof}

Combining Lemma~\ref{lem:type-completion} with
Proposition~\ref{prop:optical-residual-completion} we obtain the exact
identity
\begin{equation}\label{eq:readable-completed-height-derivative}
 D_u\mathcal F_q(w,u)\,\eta=\widehat L\eta
       +q\,\partial_\theta\bigl(B\,\mathcal A\,\eta_\theta\bigr)
       +\mathscr E_{\bx}(\mathfrak e)\,\eta,
 \qquad \mathcal A=\mathcal F_q(w,u),\quad B=\frac{E}{rQ}.
\end{equation}
Both additional terms vanish at a solution of
$(\mathcal A,\mathfrak e)=0$. We shall invert $\widehat L$ and carry
these terms as errors in the nonlinear theorem.

\subsection{Coordinates on a negative component}
\label{subsec:root-coordinate}

\paragraph{Travel coordinates.}
Let $S_\pm(r,\theta)$ be the angle at which the plus or minus
characteristic of $\widehat L$ through $(r,\theta)$ meets the left
boundary circle. Put
\begin{equation}\label{eq:fixed-strip-coordinates}
 U=h_L\circ S_+,\qquad V=h_L\circ S_-,\qquad
 \sigma=\frac{V-U}2,\qquad\varphi=\frac{U+V}2 .
\end{equation}
On the left boundary $(\sigma,\varphi)=(0,h_L)$, and by the exact
returns of $\widehat L$, on the right boundary
$(\sigma,\varphi)=(T,h_R)$. Moreover $U_\theta,V_\theta>0$ and
\[
 \det\frac{\partial(\sigma,\varphi)}{\partial(r,\theta)}
 =\widehat v\,U_\theta V_\theta>0 .
\]
For fixed $\varphi$ the function $\sigma$ increases strictly from $0$
to $T$, so $(\sigma,\varphi)$ is a diffeomorphism of the interior of
the component onto the fixed cylinder $(0,T)\times\T$. In these
coordinates the principal part of $\widehat L$, as a divergence
expression, has the density $\delta\,\diag(1,-1)$ with
$\delta=a\widehat v$: it is a wave operator, degenerate at the two
ends, where $\delta$ vanishes.

\paragraph{The root coordinate.}
At a root the speed is $\widehat v=\sqrt x\,c_0(x,\theta)$ with $c_0$
smooth and positive, so $\sigma$ is of order $x^{3/2}$. The
coordinates which are smooth at the two roots are
\[
 \xi_L=\bigl(\tfrac32\sigma\bigr)^{2/3},\qquad
 \xi_R=\bigl(\tfrac32(T-\sigma)\bigr)^{2/3},
\]
used on the halves $\sigma\le T/2$ and $\sigma\ge T/2$ of the
component (Figure~\ref{fig:root-regularization}). When the
coefficients depend on the angle, smoothness of $\xi_L$ at the root
is not obvious. To prove it, with estimates, we solve the
characteristic equation in a form in which $x^{3/2}$ is replaced by an
independent variable $\zeta$:
\begin{equation}\label{eq:fixed-strip-root-flow}
 \Theta_t=2xt\;d(xt^2,\Theta)+2\zeta t^2\,c_0(xt^2,\Theta),
 \qquad 0\le t\le1 .
\end{equation}
For $\zeta=\pm x^{3/2}$ this is the equation of the plus and minus
characteristics, parametrized by $x'=xt^2$. The inverse of the
time-one map of \eqref{eq:fixed-strip-root-flow}, followed by the
boundary chart, is one smooth function $G(x,\zeta,\theta)$, and
\[
 U=G(x,x^{3/2},\theta),\qquad V=G(x,-x^{3/2},\theta).
\]
That both families come from one function, with one boundary chart,
is essential for the cancellation in the following lemma.

\begin{lemma}[Substitution at a root]
\label{lem:rewrite-root-composition}
Let $G(x,\zeta,\theta)$ be smooth on the rectangle $0\le x\le1$,
$|\zeta|\le1$, and define
\[
 \Ev\,G=\frac{G(x,x^{3/2},\theta)+G(x,-x^{3/2},\theta)}2,\qquad
 \Od\,G=\frac{G(x,x^{3/2},\theta)-G(x,-x^{3/2},\theta)}
                         {2x^{3/2}} .
\]
Both functions extend smoothly to $x=0$, with
$\Od\,G=\partial_\zeta G(0,0,\theta)$ there, and for every integer
$s\ge0$
\begin{equation}\label{eq:rewrite-root-composition-bound}
 \|\Ev\,G\|_{C^s_{x,\theta}}\le C_s\|G\|_{C^s_{x,\zeta,\theta}},
 \qquad
 \|\Od\,G\|_{C^s_{x,\theta}}
                       \le C_s\|G\|_{C^{s+1}_{x,\zeta,\theta}} .
\end{equation}
The estimates hold also after any fixed number of differentiations
with respect to parameters.
\end{lemma}

The point of the lemma is that the half-integer power $x^{3/2}$ does
not cost half-integer many derivatives: an even function of $\zeta$
is a function of $\zeta^2=x^3$.

\begin{proof}
\step{1: $G$ even in $\zeta$}
Fix $0<x_0\le1$ and let $G_s$ be the Taylor polynomial of degree $s$
of $G$ at $(x_0,0)$ in the variables $(x,\zeta)$. Its monomials
$(x-x_0)^i\zeta^{2k}$ become $(x-x_0)^ix^{3k}$ after the
substitution. Their derivative of order $a\le s$ in $x$ at $x_0$
either vanishes or carries a factor $x_0^{3k-a+i}$ with nonnegative
exponent, and their coefficients involve at most $s$ derivatives of
$G$. For the remainder $G-G_s$, a term of the chain rule with $i$
derivatives falling on the first argument and $k$ on the second is
bounded by
\[
 C_s\|G\|_{C^s}\;
 x_0^{3(s-i-k)/2}\;x_0^{3k/2-(a-i)}
 =C_s\|G\|_{C^s}\;x_0^{3s/2-a-i/2},
\]
again with a nonnegative exponent. When the exponent is zero the
remainder still tends to zero as $x_0\to0$, by the continuity of the
derivatives of order $s$. The polynomial terms converge to the
appropriate Taylor coefficients at $x=0$. So the derivatives of the
substituted function up to order $s$ are continuous up to $x=0$ and
bounded as claimed. Derivatives in $\theta$ commute with the
construction and reduce the remaining order.

\step{2: the general case}
For a general $G$ apply Step~1 to its even part to obtain the first
bound. For the second, the function
\[
 K(x,\zeta,\theta)
 =\frac12\int_{-1}^1\partial_\zeta G(x,t\zeta,\theta)\,dt
\]
is smooth and even in $\zeta$, satisfies
$\|K\|_{C^s}\le C_s\|G\|_{C^{s+1}}$, and
$\Od\,G=K(x,x^{3/2},\theta)$; apply Step~1 to $K$. Both operations
are linear in $G$, so they commute with differentiation with respect
to parameters.
\end{proof}

Applied to the function $G$ of the flow
\eqref{eq:fixed-strip-root-flow}, the lemma gives
\[
 \varphi=\Ev\,G,\qquad
 \sigma=x^{3/2}\chi_0(x,\theta),\quad
 \chi_0(0,\theta)=\tfrac23\,c_0(0,\theta)\,h_L'(\theta)>0,\qquad
 \xi_L=x\,\bigl(\tfrac32\chi_0\bigr)^{2/3} .
\]
So $\xi_L$ is a smooth function with nonvanishing radial derivative
up to the root, and $(\xi_L,\varphi)$ are smooth coordinates there.
The right root is treated in the same way with the chart $h_R$.
Rescaling from the unit rectangle of the lemma to a collar of a root
introduces powers of the collar width and its reciprocal, which are
coefficient powers; their complete product costs one power of $M$.
The substitution $\Ev$ uses $s$ derivatives of its argument and
$\Od$ uses $s+1$, and these counts do not grow with the output order.

\begin{figure}[htbp]
\centering
\begin{tikzpicture}[x=2.5cm,y=2.0cm,font=\small,>=Stealth]
\draw[->] (0,0)--(1.45,0) node[right] {$x$};
\draw[->] (0,0)--(0,1.45);
\draw[thick,blue!65!black] plot[domain=0:1.2,samples=50] (\x,{pow(\x,1.5)});
\draw[thick,orange!85!black] (0,0)--(1.2,1.2);
\node[blue!65!black,anchor=west] at (.74,.48) {$\sigma\sim x^{3/2}$};
\node[orange!85!black,anchor=east] at (.72,1.1) {$\xi\sim x$};
\draw[dashed] (1,0)--(1,1);
\node[align=left,anchor=west] at (1.65,.85)
 {The travel coordinate $\sigma$ has zero\\derivative at the root;
  its $2/3$ power\\has a nonzero derivative.};
\end{tikzpicture}
\caption{Regularization at a turning circle; $x$ is the distance to
the root. The graph shows the leading powers with unit coefficients.
The positive coefficient of $x$ in $\xi$ is controlled uniformly at
all turning circles.}
\label{fig:root-regularization}
\end{figure}

\paragraph{The canonical form.}
Let $i\in\{L,R\}$, let $\Psi_i=(\xi_i,\varphi)$ be the coordinate map
on the corresponding half of the component, and put
\[
 \alpha_i=\frac{\delta\circ\Psi_i^{-1}}{\sqrt{\xi_i}}
 =\frac{a\,|\det D\Psi_i|}{U_\theta V_\theta}\circ\Psi_i^{-1},
 \qquad h_i=\alpha_i^{-1/2} .
\]
The second expression shows that $\alpha_i$ is smooth and positive up
to the root. In the coordinates $\Psi_i$ the principal density of
$\widehat L$ is $\alpha_i\,\diag(1,-\xi_i)$, and conjugation by
$h_i$ removes its first-order terms: for
$u=(h_iW_i)\circ\Psi_i$, the equation $\widehat Lu=f$ becomes
\begin{equation}\label{eq:fixed-strip-airy-equation}
 \begin{split}
 \bigl(\partial_{\xi_i}^2-\xi_i\,\partial_\varphi^2+V_i\bigr)W_i&=F_i,\\
 V_i&=-\alpha_i^{-1/2}\,
       \bigl(\partial_{\xi_i}^2-\xi_i\,\partial_\varphi^2\bigr)\,
       \alpha_i^{1/2},\\
 F_i&=h_i\,\Bigl(\frac{f}{|\det D\Psi_i|}\circ\Psi_i^{-1}\Bigr).
 \end{split}
\end{equation}
The principal part is the Tricomi operator
$\partial_\xi^2-\xi\,\partial_\varphi^2$, the same for all states;
only the potential $V_i$ depends on the state. The transformation of
the source preserves the pairing:
$\int uf\,dr\,d\theta=\int W_iF_i\,d\xi_i\,d\varphi$. We call
$W_i$ the \emph{canonical unknown}, $h_i$ the \emph{height factor}
and $F_i$ the \emph{canonical source}.

\paragraph{Transmission in the middle.}
The two halves meet at $\sigma=T/2$, where
$\xi_L=\xi_R=\xi_c:=(3T/4)^{2/3}$. There $\alpha_L=\alpha_R$ and
$\partial_\sigma\log\alpha_R-\partial_\sigma\log\alpha_L=\xi_c^{-3/2}$.
Continuity of $u$ and of its derivative therefore reads
\begin{equation}\label{eq:fixed-strip-cut}
 \binom{W_R}{\partial_{\xi_R}W_R}
 =\begin{pmatrix}1&0\\-1/(2\xi_c)&-1\end{pmatrix}
       \binom{W_L}{\partial_{\xi_L}W_L}
 \qquad\text{at }\xi_L=\xi_R=\xi_c .
\end{equation}
This matrix does not depend on the state.

\paragraph{Boundary data at a root.}
At a root, where the boundary chart is $h_{\rm bd}$ ($=h_L$ or $h_R$),
the value of $u$ and its inward normal flux $p_{\rm in}$ are
\begin{equation}\label{eq:fixed-strip-root-pair}
 u=(h_iW_i)(0,h_{\rm bd}),\qquad
 p_{\rm in}=\frac{h_{\rm bd}'}{h_i(0,h_{\rm bd})}\,
       \bigl(\partial_{\xi_i}W_i+(\partial_{\xi_i}\log h_i)\,W_i\bigr)
                                             (0,h_{\rm bd}) .
\end{equation}
No tangential derivative of $W_i$ appears, because
$a\varphi_r+\widehat b\varphi_\theta=0$ at the root. The identity is
exact; it retains the zero mode and the full Green pairing, not only
the principal part.

\paragraph{Bounds.}
At the reference state all these objects are explicit functions of
$r$; their bounds by coefficient powers are proved in
Lemma~\ref{lem:reference-regularized-charts}. For a general state
they are estimated in Lemma~\ref{lem:current-geometric-jets}.

\subsection{Coordinates on a positive component}
\label{subsec:positive-crossing}

\paragraph{The reference coordinate.}
On a positive component, continue the root coordinates of the two
adjacent roots across them and put $y_i=-\xi_i$, which is positive on
the positive side; here $i=L$ at the left end of the positive
component and $i=R$ at its right end. Choose a coordinate
$y\in[0,1]$, a function of $r$, equal to $y_L$ near the left end and
to $1-y_R$ near the right end, by prescribing its derivative: the
given derivatives on disjoint collars of the two ends and a positive
constant in between, interpolated with cutoffs. Similarly
interpolate the logarithms of the two height factors $h_i$, taking the
height factor equal to one in the middle. The explicit formulas are
\eqref{eq:geometry-positive-coordinate} and
\eqref{eq:geometry-positive-height}; that all these reference
quantities have coefficient bounds is
Lemma~\ref{lem:positive-reference-charts}.

\paragraph{The coordinate of a general state.}
For a general state the coordinates and height factors of the
adjacent negative components differ from the reference ones. Extend
these \emph{differences} into the positive component: the jets of the
coordinates up to order four and the jets of the logarithms of the
height factors up to order three, by Taylor polynomials multiplied
with cutoffs. The cutoffs are chosen so that the reference
coordinate is unchanged on a fixed compact interval $J_*$ of radii
which contains the support of $g_*$ in the normalization $\ell_*$ and
the support of the metric control $\chi_*$. Since the differential is
close to the reference one, the resulting map $\Psi_{\bx}$ is a
diffeomorphism onto $[0,1]\times\T$: it is proper, locally invertible,
and of degree one on the boundary. It obeys the estimates of
Proposition~\ref{prop:uniform-physical-geometry}.

\paragraph{The canonical operator.}
Let $h_{\bx}$ be the height factor in the original variables,
$\Psi_{\bx}$ the coordinate map and $J_{\bx}$ its Jacobian
determinant. For the equation $-\widehat Lu=F_{\rm ph}$ on a positive
component, the canonical operator and source are
\begin{equation}\label{eq:pp-fixed-canonical-source}
 P_{\bx}U
 =-\Bigl(\frac{h_{\bx}}{|J_{\bx}|}\,
       \widehat L\bigl[h_{\bx}\,(U\circ\Psi_{\bx})\bigr]\Bigr)
             \circ\Psi_{\bx}^{-1},\qquad
 F=\Bigl(\frac{h_{\bx}F_{\rm ph}}{|J_{\bx}|}\Bigr)
             \circ\Psi_{\bx}^{-1} .
\end{equation}
These formulas again preserve the pairing. At the reference state
\[
 P_*=-\partial_y(a\,\partial_y)-c\,\partial_\varphi^2+V,
 \qquad a=h_*^2\,w_*\,y_r,\quad c=\frac{h_*^2\,w_*'}{r\,y_r},
\]
where in this formula $a$ and $c$ denote the coefficients of $P_*$.
Near a root $a=1$ and $c=y_i$ exactly: $P_*$ is the continuation
$-\partial_{y_i}^2-y_i\,\partial_\varphi^2$ of the Tricomi operator.
On the whole component, with
\[
 \bd=y\,(1-y),
\]
we have
\begin{equation}\label{eq:positive-coefficients}
 a\ge\mathcal E_m^{-C},\qquad c\ge\mathcal E_m^{-C}\,\bd ,
\end{equation}
and all derivatives of the coefficients of fixed order are bounded by
coefficient powers
(Lemma~\ref{lem:positive-reference-charts}); away from the roots
this comes from the lower bounds for $w_*'/w_*$ on the positive parts
of the models and on the joins.

\paragraph{The difference vanishes to third order at the roots.}
Since the jets specified above match the principal densities of
$P_{\bx}$ and $P_*$ to second order at both roots,
\begin{equation}\label{eq:readable-positive-difference}
 P_{\bx}-P_*=-\partial_a\bigl(\bd^{\,3}\,e_{\bx}^{ab}\,\partial_b\bigr)
                           +\delta V_{\bx} ,
\end{equation}
with smooth coefficients $e^{ab}_{\bx}$, where $a,b$ are indices
running over $y,\varphi$, and a potential $\delta V_{\bx}$ on which no
boundary condition is imposed. Division by $\bd^3$ has a fixed cost:
if $F$, $F'$ and $F''$ vanish at $0$, then
\[
 \frac{F(y)}{y^3}=\frac12\int_0^1(1-t)^2\,F'''(ty)\,dt ,
\]
each further derivative differentiates the integrand, and the factor
$t^{-1/2}$ produced by rescaling in $L^2$ is integrable. This is the
content of part~(iii) of
Proposition~\ref{prop:uniform-physical-geometry}
(Lemma~\ref{lem:current-geometric-jets}, Step~4). The
cubic vanishing is what will allow a second-order perturbation to be
controlled without boundary conditions
(Section~\ref{subsec:cubic-estimate}).

\subsection{Proof of the uniform geometric calculus}
\label{subsec:geometry-proofs}

We prove Proposition~\ref{prop:uniform-physical-geometry}. The bounds
at the reference state come first; they involve functions of $r$ only
(Lemmas~\ref{lem:reference-regularized-charts}--\ref{lem:geometric-weight-conversion}).
The bounds for a general state follow from them and from the lemmas
on flows
(Lemmas~\ref{lem:uniform-characteristic-flow-constants}
and~\ref{lem:current-geometric-jets}). Positivity of a function is
never used as a bound for its reciprocal: every reciprocal that
occurs is bounded explicitly.

\begin{lemma}[Root charts of the reference profile]
\label{lem:reference-regularized-charts}
\label{lem:reference-height-source-factors}
On each half of the negative component of cell $n$, the root
coordinate of the reference profile is the function of $r$
\begin{equation}\label{eq:geometry-regular-root-coordinate}
 \xi=x\,\Bigl[\tfrac32\int_0^1\sqrt t\;c_i(xt)\,dt\Bigr]^{2/3},
\end{equation}
where $x$ is the distance from the root of that half, and
\begin{equation}\label{eq:geometry-root-inverse}
 \mathcal E_n^{-C}\le|\xi_r|\le\mathcal E_n^{C},\qquad
 \|\xi\|_{C^s}+\|r(\xi)\|_{C^s}\le C_s\,\mathcal E_n^{C_s}.
\end{equation}
Moreover
\begin{equation}\label{eq:geometry-root-density}
 \alpha_i=w_*\,|\xi_r|,\qquad h_i=\alpha_i^{-1/2},
\end{equation}
and the functions $\alpha_i^{\pm1}$, $h_i^{\pm1}$, $|\xi_r|^{\pm1}$,
the potential $V_i=-\alpha_i^{-1/2}\partial_\xi^2\alpha_i^{1/2}$, the
source factors $h_i/|\xi_r|$ and $|\xi_r|/h_i$, and
$\partial_\xi\log h_i$ have derivatives of order $s$ bounded by
$C_s\mathcal E_n^{C_s}$. These bounds persist when the half is
rescaled to unit length.
\end{lemma}

\begin{proof}
At the reference state the charts are the identity, so
$\varphi=\theta$ and, on each half, $\xi=(\tfrac32\sigma_i)^{2/3}$,
where $\sigma_i=\int v_0\,dr$ is the travel angle from the root of
that half.

\step{1: near the root}
Since $v_0=\sqrt x\,c_i(x)$, the substitution $x'=xt$ gives
$\sigma_i=x^{3/2}\int_0^1\sqrt t\,c_i(xt)\,dt$, which is
\eqref{eq:geometry-regular-root-coordinate}. The bracket is smooth
and positive, and equals $c_i(0)$ at $x=0$; so
$|\xi_r|=c_i(0)^{2/3}$ at the root. On the collar
$|x|<8\delta_n^{\rm root}$ the bracket lies between
$\tfrac12c_i(0)$ and $2c_i(0)$, and by
\eqref{eq:geometry-normalized-root-jets} its logarithmic derivative
times $|x|$ is at most $C\epsilon_{\rm root}A_n^{15}$. Hence
$|\xi_r|$ is comparable to $c_i(0)^{2/3}$ on the collar, and the
derivatives of $\xi$ have coefficient bounds there.

\step{2: away from the collars}
Differentiating $\xi^{3/2}=\tfrac32\sigma_i$ gives
$|\xi_r|=v_0/\sqrt\xi$, so we need upper and lower bounds for $v_0$
and $\xi$. In the model coordinate the collar of the left root ends
at a distance comparable to $A_n^{19}$ from $X_L$, and that of the
right root at a distance comparable to $A_n^{15}$ from $X_R$, because
$r_X$ is comparable to $\mathfrak r_nA_n$ and $\mathfrak r_nA_n^5$
there. The function $\rho$ has one maximum on the component and
vanishes only at its two ends, where $|\rho_X|$ is comparable to
$A_n$ and $A_n^9$. Hence $\rho\ge cA_n^{24}$ outside the collars, and
\[
 c\,\mathfrak r_n^{-1}A_n^{11}\le v_0=\frac{\sqrt\rho}{r_X}
 \le C\,\mathfrak r_n^{-1}A_n^{-9/2}.
\]
By Step~1, $\xi\ge cA_n^{20}$ at the end of its collar, hence beyond
it, and $\xi\le\xi_c=(3T_n/4)^{2/3}\asymp A_n^{1/3}$. So $|\xi_r|$ is
bounded above and below by powers of $A_n$. The higher derivatives
follow by differentiating $\xi_r=\pm v_0\,\xi^{-1/2}$, using the lower
bounds for $r_X$, $\rho$ and $\xi$ on this region.

\step{3: the inverse}
From $r(\xi(r))=r$ we get $r_\xi=1/\xi_r$, and every higher
derivative of $r(\xi)$ is a polynomial in lower ones and in
derivatives of $\xi$, divided by $\xi_r$. This proves
\eqref{eq:geometry-root-inverse}. The length $\xi_c$ of a half and
its reciprocal are powers of $A_n$, so rescaling to unit length costs
a fixed power of $A_n$ for each derivative.

\step{4: the factors}
For the reference operator $a=w_*$, $\widehat b=0$,
$\delta=w_*v_0$, $U_\theta=V_\theta=1$ and
$|\det D\Psi_i|=|\xi_r|$. Hence
$\alpha_i=\delta/\sqrt\xi=w_*|\xi_r|$, which is
\eqref{eq:geometry-root-density}. By
\eqref{eq:geometry-profile-jets} and \eqref{eq:geometry-root-inverse}
the functions $\alpha_i^{\pm1}$ have coefficient bounds, and the
remaining quantities are obtained from $\alpha_i$ and $\xi_r$ by
products, powers and differentiation in $\xi$. All of them are
functions of $r$, so no angular frequency is involved.
\end{proof}

\begin{lemma}[The reference coordinate on a positive component]
\label{lem:positive-reference-charts}
\label{lem:positive-reference-coefficient-bounds}
Let $(r_1,r_2)$ be a positive component with index $n$, let
$y_L,y_R$ be the continued root coordinates at its two ends, and let
$\eta_L,\eta_R$ be cutoffs with disjoint supports in the collars of
the two ends, equal to one on smaller collars. Define
\begin{equation}\label{eq:geometry-positive-coordinate}
 \begin{split}
 y_r&=\eta_L\,(y_L)_r-\eta_R\,(y_R)_r+\mu_0\,(1-\eta_L-\eta_R),
 \qquad y(r)=\int_{r_1}^ry_r\,dr,\\
 \mu_0&=\frac{1-\int\eta_L\,(y_L)_r\,dr+\int\eta_R\,(y_R)_r\,dr}
            {\int(1-\eta_L-\eta_R)\,dr},
 \end{split}
\end{equation}
and
\begin{equation}\label{eq:geometry-positive-height}
 \log h_*=\eta_L\log h_L+\eta_R\log h_R,
 \qquad h_i=\bigl(w_*\,|(y_i)_r|\bigr)^{-1/2}.
\end{equation}
Then $y$ is a diffeomorphism of $[r_1,r_2]$ onto $[0,1]$ with
$y=y_L$ near $r_1$ and $y=1-y_R$ near $r_2$, and $y_r^{\pm1}$,
$h_*^{\pm1}$, the inverse function $r(y)$ and the source factors
$h_*/y_r$, $y_r/h_*$ have derivatives of order $s$ bounded by
$C_s\mathcal E_n^{C_s}$. The coefficients of the reference operator
$P_*=-\partial_y(a\,\partial_y)-c\,\partial_\varphi^2+V$ are
\begin{equation}\label{eq:geometry-positive-coefficients}
 a=h_*^2\,w_*\,y_r,\qquad
 c=\frac{h_*^2\,w_*'}{r\,y_r},\qquad
 V=-\frac{h_*}{y_r}\,\bigl(w_*\,h_{*,r}\bigr)_r .
\end{equation}
Near the ends $a=1$ and $c=y_i$; on the whole component
\eqref{eq:positive-coefficients} holds, and the derivatives of order
$s$ of $a$, $c$, $V$ and $a^{-1}$ are bounded by
$C_s\mathcal E_n^{C_s}$. On the component which contains the interval
$J_*$, the collars can be chosen disjoint from $J_*$.
\end{lemma}

\begin{proof}
\step{1: the collars}
The continued coordinates $y_i=-\xi_i$ have positive inward
derivatives with the bounds of
Lemma~\ref{lem:reference-regularized-charts}, since
\eqref{eq:geometry-normalized-root-jets} holds on both sides of the
root. The increment of $y_i$ over its collar is at most
$8\delta^{\rm root}\sup|\xi_r|=O(A_n^{18})$. So, after
$\epsilon_{\rm root}$ is decreased if necessary and finitely many
central collars are chosen by hand, the two collars together account
for less than $\tfrac14$ of the total increment $1$. The interval
$J_*$ has positive distance from the roots.

\step{2: the coordinate}
Note that $(y_R)_r<0$, since $y_R$ increases inwards from the right
end. The numerator of $\mu_0$ lies between $\tfrac34$ and $1$, and its
denominator is comparable to $r_2-r_1$; so $\mu_0$ and $\mu_0^{-1}$
are coefficient powers. The function $y_r$ is a convex combination of
positive functions with coefficient bounds above and below, and its
integral over the component is $1$ by the choice of $\mu_0$. Hence
$y=y_L$ near $r_1$ and $y=1-y_R$ near $r_2$, and the inverse function
is estimated as in Step~3 of the preceding proof.

\step{3: the height factor}
The logarithms $\log h_i$ are bounded by $C(n^{5/2}+n^2)$ by
\eqref{eq:geometry-profile-jets}, and their derivatives of positive
order have coefficient bounds. Differentiating
\eqref{eq:geometry-positive-height} and exponentiating gives the
bounds for $h_*^{\pm1}$.

\step{4: the coefficients}
The formulas \eqref{eq:geometry-positive-coefficients} follow by
inserting $u=h_*\,U(y(r))\cos(j\theta)$ into
$-(h_*/y_r)\,\mathscr L_{w_*}u$; the terms with $U_y$ which do not
belong to $(aU_y)_y$ cancel. In a collar
$h_*=(w_*\,|(y_i)_r|)^{-1/2}$ gives $a=1$. Differentiating
\eqref{eq:geometry-regular-root-coordinate} gives
$\xi\,\xi_r^2=x\,c_i(x)^2=-w_*'/(rw_*)$ on both sides of the root, so
$w_*'/(rw_*)=y_i\,(y_i)_r^2$ and $c=y_i$. The upper bounds and the
lower bound for $a$ follow from Steps~2 and~3.

\step{5: the lower bound for $c$}
It suffices to show
\begin{equation}\label{eq:geometry-wprime-distance}
 w_*'(r)\ge\mathcal E_n^{-C}\,\min(r-r_1,\,r_2-r),
\end{equation}
because $y(1-y)\le\mathcal E_n^{C}\min(r-r_1,r_2-r)$ by the bounds
for $y_r$. On a short positive cap
$-\rho\ge cA_n\min(|X-X_L|,1)$, and on a long positive cap
$-\rho\ge cA_n^9\min(|X-X_R|,1)$; since
$w_*'=-\lambda_n\rho\,r/r_X$, the bounds of Lemma~\ref{lem:cell} give
\eqref{eq:geometry-wprime-distance} on the caps. On a join
$w_*'/w_*\ge2/r$ by \eqref{eq:join-slope}, and on the connectors of
Lemma~\ref{lem:connector} $w_*'/w_*\ge3/(4r)$; with the lower bound
for $w_*$ this gives \eqref{eq:geometry-wprime-distance} there, since
the distance to the ends is bounded.
\end{proof}

\begin{lemma}[Conversion of weights and derivatives]
\label{lem:geometric-weight-conversion}
On cell $n$ and its neighbours, the continuous weights $H$, $E_0$,
$M$ and the discrete weights $H_n$, $\mathcal E_n$, $M_n$ differ by
factors bounded by $\mathcal E_n^{C}$. Changes between the variables
$r$, $\varrho$, the root coordinates and the coordinates $y$ transfer
derivatives of total order $s$ to derivatives of total order at most
$s$, with coefficients bounded by $E_0^{C_s}$; so do the source
factors and cutoffs. A restriction to a circle costs one radial
derivative.
\end{lemma}

\begin{proof}
By \eqref{eq:geometry-endpoint-distance},
$x_\partial=n+O(1)$ on cell $n$, so
\[
 \log\frac H{H_n}=O(n^2),\qquad
 \log\frac{E_0}{\mathcal E_n}=O\bigl(n^{3/2}\log(e+n)\bigr),\qquad
 \log\frac M{M_n}=O\bigl(n^{3/2}\log^2(e+n)\bigr),
\]
and the same bounds hold for the ratios of the weights of
neighbouring cells; all of them are $O(\log\mathcal E_n)$. In the
variable $\varrho$,
\[
 |\Dr\log M|\le C\,(1+x_\partial)^{-1/2}\log^2(e+x_\partial),
\]
which is bounded together with its derivatives, and the same is true
of $\log H$, $\log E_0$ and $\log\ell$; so the weights are uniformly
comparable on unit intervals of $\varrho$, and $\ell^{\pm1}$ is
bounded by a power of $E_0$. The chain rule expresses derivatives of
total order $s$ in one set of variables through derivatives of order
at most $s$ in another, with coefficients which are polynomials in the
derivatives of the coordinate functions and in reciprocals of their
first derivatives; these are bounded by
Lemmas~\ref{lem:reference-regularized-charts}
and~\ref{lem:positive-reference-charts}. The inequality
$E_0^{C_s}\le C_sM$ is applied to the complete product. The trace
inequality on an interval of unit length costs one derivative, and
the rescaling to unit length costs coefficient powers.
\end{proof}

\begin{lemma}[Flows at a general state]
\label{lem:uniform-characteristic-flow-constants}
On the neighbourhood of
Lemma~\ref{lem:common-geometric-neighborhood}, after $b_0$ is
increased and $\epsilon_0$ decreased, the characteristic speeds
$c_\pm=d\pm v$ of $\bar L$ on a negative component $m$ satisfy, away
from the collars of the roots,
\begin{equation}\label{eq:geometry-flow-low-bound}
 \int\|\partial_\theta c_\pm(r,\cdot)\|_{H^s_\theta}\,dr
 \le C_s\,\mathcal E_m^{C_s}\,H_m^\beta\;\mathfrak g_{s+d_*,m},
\end{equation}
and the right sides of the regularized equation
\eqref{eq:fixed-strip-root-flow} in the collars, and the corrected
speeds $\widehat c_\pm$ of
Proposition~\ref{prop:optical-residual-completion}, satisfy the same
bound. Consequently the angular derivatives of all characteristic
flows and of their inverses lie in $[\tfrac12,2]$, and the constants
of Lemma~\ref{lem:circle-flow} for these flows are uniform in $m$.
\end{lemma}

\begin{proof}
\step{1: the speeds}
By \eqref{eq:completed-characteristics},
\[
 d=\frac{E\,(p-k_\theta)}{rQ},\qquad
 v=\sqrt{-\frac{w'}{rw}}\;\frac{w\sqrt D}{Q},
 \qquad -\frac{w'}{rw}=v_0^2-\frac{f'}r,
\]
where $w=w_*e^f$. At the reference state $d=0$ and $v=v_0$. The
quantities $d$ and $w\sqrt D/Q-1$ are coefficients in the sense of
Step~3 of the proof of Lemma~\ref{prop:full-physical-tame}, and they
vanish at the reference state. Outside the supports of the controls
$f'=0$; on the support of $\chi_n$ the ratio $f'/(rv_0^2)$ is bounded
by $C|\kappa_n|A_n^{-29}$, with derivatives bounded by coefficient
powers times $|\kappa_n|$
(Lemma~\ref{lem:quantitative-optical-controls}). So $v/v_0-1$ is of
the same kind. Only the deformation depends on $\theta$, and the
length of the component is bounded; this proves
\eqref{eq:geometry-flow-low-bound}. In a collar write
$v=\sqrt x\,c_0(x,\theta)$ with
$c_0=c_i(x)\cdot(v/v_0)$; by Lemma~\ref{lem:quantitative-root-collars}
the function $c_0$ and its reciprocal have coefficient bounds, and
$c_0-c_i$ and $d$ are bounded as in
\eqref{eq:geometry-flow-low-bound}. The equation
\eqref{eq:fixed-strip-root-flow} involves $d$ and $c_0$ at the point
$xt^2$ and the bounded factors $x$ and $\zeta$, so its right side has
the same bounds.

\step{2: the Jacobians}
Rescale the radial variable on each part of the component to the unit
interval; this multiplies the speeds by lengths, which are coefficient
powers. By \eqref{eq:geometry-smallness-absorption}, the right side of
\eqref{eq:geometry-flow-low-bound} for $s\le4$ is as small as we wish
on the neighbourhood, uniformly in $m$. Lemma~\ref{lem:circle-flow}
then gives the bounds for the angular derivatives of the flows; the
derivative of an inverse flow is the reciprocal of that of the flow.
The large speed $v_0$ of the reference profile does not enter,
because it does not depend on $\theta$.

\step{3: the corrected flow}
The corrected minus flow is $\widehat\Phi_r=\Phi_r\circ J_r$ with
$J_r=\id+\chi(r)\Delta$, and its speed is
\eqref{eq:readable-optical-velocity}. At the reference state
$\Delta=0$. By the definition of $\Delta$ in the proof of
Proposition~\ref{prop:optical-residual-completion} and
Lemma~\ref{lem:one-high-factor}(iv), $\Delta$ is estimated by the
return residual, hence by the right side of
\eqref{eq:geometry-flow-low-bound}; the cutoff $\chi$ has
coefficient bounds, and on the support of $\chi'$ the speed $v$ is
bounded below by a coefficient power, so the correction is smaller
than $v/2$. It vanishes near both roots, so the functions $c_0$ are
unchanged. Thus Steps~1 and~2 apply to the corrected flows.
\end{proof}

\begin{lemma}[Coordinates and factors at a general state]
\label{lem:current-geometric-jets}
On the same neighbourhood, with a fixed $d_*$, every coordinate,
inverse coordinate, Jacobian, height factor, source factor, potential
and divided coefficient of
Sections~\ref{subsec:return-correction}--\ref{subsec:positive-crossing}
satisfies \eqref{eq:geometry-current-jet-bound}, and its derivatives
with respect to the state satisfy the corresponding estimates of the
shape \eqref{eq:tame-derivative-shape}.
\end{lemma}

\begin{proof}
We follow the order in which the objects were constructed. No norm of
high order is assumed small.

\step{1: flows and root coordinates}
Away from the collars, the last statement of
Lemma~\ref{lem:circle-flow} bounds the difference between the flows
of a general state and of the reference state by the difference of
the speeds, that is, by the right side of
\eqref{eq:geometry-flow-low-bound}. In a collar the same lemma
applies to \eqref{eq:fixed-strip-root-flow}, whose right side is
smooth in the independent variables $(x,\zeta,\theta)$; so the
function $G(x,\zeta,\theta)$ of Section~\ref{subsec:root-coordinate},
minus its reference value, satisfies
\eqref{eq:geometry-current-jet-bound} in all three variables.
Lemma~\ref{lem:rewrite-root-composition} then gives the bound for
$\varphi=\Ev\,G$, and for $\chi_0=\Od\,G$ with one additional
derivative; and $\xi_L=x\,(\tfrac32\chi_0)^{2/3}$ is a smooth function
of $\chi_0$, which is bounded below. The charts $h_R$ and the residual
$\mathfrak e$ are obtained from $h_L$ and the return map by
compositions, inversions and parity projections, to which
Lemma~\ref{lem:one-high-factor}(iv) applies. The corrected flows were
treated in Lemma~\ref{lem:uniform-characteristic-flow-constants}.

\step{2: inverse coordinates, height factors and sources}
At low order, Step~1 and \eqref{eq:geometry-root-inverse} give, for
the coordinate maps rescaled to unit size,
\[
 \bigl\|(D\Psi-D\Psi_*)\,(D\Psi_*)^{-1}\bigr\|_{C^0}
 \le C\,\epsilon_0\,H_m^\beta\,\mathcal E_m^{C}<\tfrac14
\]
on the neighbourhood, by \eqref{eq:geometry-smallness-absorption}. So
$D\Psi$ is invertible with bounds comparable to the reference. The
higher derivatives of $\Psi^{-1}$ follow by differentiating
$\Psi^{-1}\circ\Psi=\id$; subtracting the reference identity and
using the product estimate keeps the factor of the difference. For
the density use the identity
\[
 \alpha_i=\frac{a\,|\det D\Psi_i|}{U_\theta V_\theta}\circ\Psi_i^{-1}
\]
of Section~\ref{subsec:root-coordinate}: the Jacobian of the
regularized coordinates does not vanish, and $U_\theta V_\theta$ is
bounded below by
Lemma~\ref{lem:uniform-characteristic-flow-constants}. Powers,
derivatives and the formulas \eqref{eq:fixed-strip-airy-equation}
then give the height factor, the potential and the source factor. The
reciprocal of the Jacobian of the unregularized coordinates
$(\sigma,\varphi)$, which vanishes at the roots, is never used.

\step{3: positive components}
The coordinate and the height factor of a general state on a positive
component are the reference ones plus Taylor polynomials, of degree
four and three, of the differences at the two ends, multiplied by
fixed cutoffs. Restriction to the boundary, the polynomials and the
cutoffs preserve the bounds of Steps~1 and~2. The differential stays
close to the reference one as in Step~2, so the Jacobian is positive,
and the inverse coordinate and the source factor are estimated in the
same way.

\step{4: division by $\bd^3$}
The differences $F$ of the principal densities vanish with their
first two inward derivatives at both ends. Near the left end
\[
 \frac{F(y)}{y^3}=\frac12\int_0^1(1-t)^2\,F'''(ty)\,dt,
\]
which costs three derivatives and keeps the factor $H_m^\beta$; near
the right end divide by $(1-y)^3$ in the same way; in the middle
$\bd^{-1}$ has coefficient bounds. This gives the bound for
$e^{ab}_{\bx}=F^{ab}/\bd^3$.

\step{5: the loss and the parameters}
The characteristic speeds involve two derivatives of the height; the
divided difference at a root one more; the boundary polynomials four;
the principal density one derivative of the coordinates; the potential
two more; the division three; and the Sobolev embedding and the
traces a fixed number. Their total is $d_*$, which does not depend on
$s$. Derivatives with respect to the state are estimated by
differentiating the same finite formulas and the equations of the
flows: every variation of a flow solves a linear equation with the
same integrating factor, which is bounded by Step~2 of the preceding
proof. The height enters through $Z=qu$, whose first derivative is
$\dot q\,u+q\,\dot u$, whose second derivative is $q_1u_2+q_2u_1$, and
whose higher derivatives vanish. Hence the seminorms
\eqref{def:joint-displacement-seminorms}, which contain the unscaled
$u$, control all derivatives, also at $q=0$. Finally
Lemma~\ref{lem:geometric-weight-conversion} absorbs each complete
product of reference functions into one power of $M$.
\end{proof}

\begin{proof}[Proof of Proposition~\ref{prop:uniform-physical-geometry}]
Part~(i) is Lemmas~\ref{lem:reference-regularized-charts}
and~\ref{lem:positive-reference-charts}, with
Lemmas~\ref{lem:quantitative-root-collars}
and~\ref{lem:quantitative-optical-controls} for the collars, the
controls and the cutoffs. Parts~(ii) and~(iii) are
Lemmas~\ref{lem:uniform-characteristic-flow-constants}
and~\ref{lem:current-geometric-jets}; the comparability of the
differentials was shown in Step~2 of the last proof. All collars and
cutoffs were chosen before any order of differentiation; only the
constants depend on the order.
\end{proof}

\subsection{The fixed transmission domain}
\label{subsec:transmission-domain}

We have represented a height on the annulus by canonical unknowns on a
collection of cylinders: two half-cylinders
$[0,\xi_c]\times\T$ for each negative component and one cylinder
$[0,1]\times\T$ for each positive component. The cylinders do not
depend on the state. Neither do the conditions that link them:
\begin{itemize}
\item at the middle of a negative component, the condition
 \eqref{eq:fixed-strip-cut};
\item at a root, equality of the canonical values and of the inward
 derivatives with opposite signs. Indeed, the coordinates and the
 height factors of the two adjacent components have the same value
 and the same first derivative at the root, so
 \eqref{eq:fixed-strip-root-pair} takes this form.
\end{itemize}
We call the set of canonical unknowns satisfying these conditions the
\emph{transmission domain}. A smooth element of it which solves the
equation corresponds to a function which is smooth in the original
variables across the roots: equality of the values gives equality of
all tangential derivatives; equality of the fluxes gives equality of
the first normal derivative; and the equation can be solved for
$a\,u_{rr}$, because $a>0$, so that repeated differentiation matches
all higher normal derivatives.

This is the result of the section: all the operators
$\widehat L_{\bx}$ act, in canonical form, on one domain. It is what
makes it possible to regard them as perturbations of the reference
operator in the next section.
